% !TEX root = ../../../main.tex
% Block card: preprocessing.sliding_windows -- "Sliding windows" (Part I, the blocks chapter).
\cardfield{What it does}
Cuts the span into equal, non-overlapping windows and marks each one as training,
validation or test, keeping whole stretches of time together.

\cardfield{When to use it}
At the start of any chain that will train or test something on windows: a classifier, a
clustering you want to check on unseen time, window images or MiniRocket. The windows
carry no measurements yet; put a feature block after it. To pool windows from several
recordings and channels, use \blockref{preprocessing.window_pool} instead. The Library's
window sets explain the same rules from the page side: \apppage{library.window-sets}.

\begin{plainterms}[Why a blocked split]
A fungal recording changes slowly. Two windows next to each other in time are near-copies.
If you shuffle windows at random into ``train'' and ``test'', the test set is full of
near-twins of training windows, and the model looks far better than it is: like marking a
student's exam with questions copied from their homework. A \emph{blocked split} cuts the
recording into a few long blocks of time and gives each block, whole, to one role. The test
windows then come from a stretch of time the model never saw.
\end{plainterms}

\cardfield{What goes in, what comes out}
In goes a span of one channel. Out comes a set of window positions, each with its role
(train, validation or test) attached. The windows touch end to end unless you ask for a gap.

\cardfield{Settings that matter}
\emph{window\_s} is the window length in seconds (default ten minutes). \emph{n\_blocks}
is how many blocks of time the span is cut into; \emph{holdout\_frac} and
\emph{validation\_frac} say what share of those blocks becomes test and validation. The
last blocks in time are the ones held out.

\cardfield{Watch out for}
The share of blocks is rounded: with five blocks a validation share of 0.1 rounds to no
validation block at all. There is no empty gap at the border between a training block and
the test block, so the last training window and the first test window can sit side by
side. A random split is refused unless you tick \emph{allow\_random\_split}, and you
should not.
